\documentclass{article}

\usepackage{amsmath,amssymb}
\usepackage{xurl}
\usepackage[hidelinks]{hyperref}
\setlength{\emergencystretch}{2em}

\input{./command}
\providecommand{\leanid}[1]{\path{#1}}

\title{The nonempty-family condition in Wolf's Radon--Nikodym theorem}
\author{}
\date{2026-08-23}

\begin{document}
\maketitle
\gapnote{local-correction}{resolved}

\section{Source statement}

Wolf's Radon--Nikodym theorem for quantum instruments states that, if
\(\{T_i\}\) is a family of completely positive linear maps
\[
  T_i,T:\mathcal M_{d'}\longrightarrow\mathcal M_d
\]
satisfying \(\sum_i T_i=T\), and if a Stinespring representation of the total
map is supplied by
\[
  V:\mathbb C^d\longrightarrow\mathbb C^{d'}\otimes\mathbb C^r,
  \qquad
  T(A)=V^\dagger(A\otimes \mathbf 1_r)V,
\]
then there are positive operators \(P_i\in\mathcal M_r\) such that
\[
  \sum_i P_i=\mathbf 1_r,
  \qquad
  T_i(A)=V^\dagger(A\otimes P_i)V.
\]
The local source is
\path{Notes/WolfNoteTexSource/ch02_representations.tex}:402--407, titled
``Radon--Nikodym for quantum instruments'' in
\cite{Wolf2012Quantum}.

\section{Local boundary}

The source-facing Lean theorem
\leanid{IsKrausCP.radon_nikodym_of_stinespring} makes the finite index type
nonempty. This is not a strengthening of a mathematical hypothesis in the
substantive case; it is the finite-family boundary needed for the identity
\[
  \sum_i P_i=\mathbf 1_r
\]
to be meaningful as a resolution of the identity by instrument components.  If
the index type were empty, then the hypotheses \(\sum_i T_i=T\) force
\(T=0\).  One may still supply the zero Stinespring matrix
\(V=0\), but the conclusion would ask
\[
  0=\mathbf 1_r
\]
in \(\mathcal M_r\), which is false for a nonzero dilation dimension.

Thus the formal theorem states the source theorem in the nonempty instrument
case, with its rectangular Heisenberg dimensions made explicit. The older
\leanid{IsCPMap.radon_nikodym_of_stinespring} theorem is retained as the
square-algebra specialization.

The formal proof follows the Kraus-freedom route. Choose rectangular Kraus
families for the component maps, combine them into one labelled family, pad one
component by zero Kraus operators, and compare that family with the adjoints
of the ancilla blocks of the supplied Stinespring matrix. Rectangular Kraus
freedom gives an isometry whose fibre Gram matrices are the
operators \(P_i\). No Douglas-factorization argument is used.

\section{Elimination plan}

No mathematical proof gap is expected for the nonempty case.  An empty-index
variant would need a different conclusion, for example a vacuous statement with
zero dilation space or no identity resolution.  Such a statement is not the
instrument theorem used in the development.

Verdict: local correction of a well-posedness boundary in the printed finite-family
statement.

\bibliographystyle{alpha}
\bibliography{references}

\end{document}
